\section{The acx-cert-v1 Certificate Schema and Its Verification}\label{app:certs}

Every solution and every stable-AC equivalence claimed in this paper is exported as an
acx-cert-v1 JSON certificate and checked by two independently written verifiers before it
is reported.

\subsection{Schema}

A certificate is a JSON object with fields: \texttt{certificate\_version} (schema version
string); \texttt{name}; \texttt{claim} (a human-readable statement of what is being
certified); \texttt{start} and \texttt{end} (each an object \texttt{\{n\_gen, relators\}});
\texttt{end\_is\_trivial} (boolean, true iff \texttt{end} must equal the trivial
presentation on \texttt{n\_gen} generators); \texttt{steps} (an ordered list of length
$T$); \texttt{states} (the list of intermediate presentations, length $T+1$, with
\texttt{states[0]==start} and \texttt{states[-1]==end}); and \texttt{meta} (free-form
provenance --- source paper or line numbers, generator naming, and similar bookkeeping).

\subsection{Step types}

Six step types are implemented; the field names below are exact, not paraphrased:

\begin{itemize}
  \item \textbf{concat} \texttt{\{i, j, sign\}} --- $r_i \to$ freely-reduced$(r_i \cdot
    r_j^{\text{sign}})$ (AC1; \texttt{sign=-1} concatenates with $r_j^{-1}$).
  \item \textbf{conjugate} \texttt{\{i, g\}} --- $r_i \to$ freely-reduced$(g\cdot r_i\cdot
    g^{-1})$ (AC3).
  \item \textbf{stabilize} \texttt{\{z, w\}} --- appends generator $z=$ \texttt{n\_gen}$+1$
    and the cyclically-reduced relator $z\cdot w^{-1}$ (AC4 plus realizing $z=w$).
  \item \textbf{eliminate} \texttt{\{gen, ri, ...\}} --- Lemma-11 removal of \texttt{gen}
    using its exactly-once occurrence in relator \texttt{ri}; generators above \texttt{gen}
    are renumbered down by one (the packaged AC5).
  \item \textbf{relabel} \texttt{\{perm, invert\}} --- a bijective, sign-respecting
    permutation of the generator set: an AC-equivariant renaming, not a move on relators.
  \item \textbf{substitution} \texttt{\{ci, a, b, i, j, c\_inv\}} --- the composite move used
    by the base greedy solver (\ref{sec:methods}); verified by checking that the resulting
    state lies in the recomputed neighbor set of the preceding state, rather than by
    replaying finer sub-steps.
\end{itemize}

There is no separate step type for a bare AC2 inversion of a whole relator. A full-relator
sign flip is instead absorbed by the canonical-equality quotient the verifiers already
apply (reorder relators, cyclically rotate, invert) --- see the mutation-sensitivity note
below, where this distinction mattered in practice.

\subsection{Verification levels}

\begin{itemize}
  \item \textbf{L1 (replay).} Recompute each step's post-state from its recorded parameters
    and the preceding state; it must equal the next recorded state exactly.
  \item \textbf{L2 (preconditions).} \texttt{eliminate}'s generator must occur exactly once
    in the cited relator; \texttt{relabel}'s permutation must be bijective and
    sign-consistent; every letter must lie in the valid signed-generator range for that
    state's \texttt{n\_gen}.
  \item \textbf{L3 (global invariants).} Every state must be balanced (\texttt{n\_gen}
    equal to the relator count, no empty relator); the abelianization matrix's determinant
    must have absolute value $1$ at every state, computed by exact-integer Bareiss
    elimination (no floating point); if \texttt{end\_is\_trivial} is claimed, the end state
    must equal $\langle x_1,\dots,x_k\mid x_1,\dots,x_k\rangle$ up to relator reordering.
\end{itemize}

\subsection{Two independent verifiers}

Every certificate is checked by the engine author's verifier and by a second verifier
authored black-box from the schema and the math above alone, with no access to the engine's
source and its own free/cyclic reduction, canonicalization, integer-determinant, and
step-replay logic. Across the campaign the two verifiers together performed $20{,}947$
checks with $0$ failures. All five exported chain certificates pass both
(Table~\ref{tab:certs}): \texttt{appendixF\_P25\_to\_AK3} (53 steps, the literature's
$\mathrm{P25}\to\mathrm{AK}(3)$ replay of \ref{app:misprint}), \texttt{laneF\_F\_to\_AK3}
(21 steps, the path from $F$ --- a 2-generator presentation AC-equivalent to AK(3) --- to AK(3)
itself), and three small stable-solver certificates recording in-search destabilizations
(3, 1, and 1 steps respectively).

\begin{table}[tbp]
\centering
\caption{The five exported acx-cert-v1 chain certificates, each passing both the engine
author's verifier and the independently authored black-box verifier.}
\label{tab:certs}
\small
\setlength{\tabcolsep}{3pt}%
\begin{tabular}{@{}p{0.27\textwidth}p{0.55\textwidth}rr@{}}
\toprule
Certificate & Claim & Steps & States \\
\midrule
appendixF\_P25\_to\_AK3 & P25 is AC-equivalent to AK(3) via the 53 printed Appendix-F h-moves (arXiv:2408.15332v2). & 53 & 54 \\
laneB\_M3corr\_hero8\_500 & Stable-AC trivialization of M3corr\_hero8 found by StableSolver. & 3 & 4 \\
laneB\_ms0 & Stable-AC trivialization of ms0\_plain found by StableSolver. & 1 & 2 \\
laneB\_ms0\_stab & Stable-AC trivialization of ms0\_stab found by StableSolver. & 1 & 2 \\
laneF\_F\_to\_AK3 & F = \textless{}x,y | YYxyXX, YYYXXyx\textgreater{} is plain-AC-equivalent to AK(3) (up to signed relabeling); explicit substitution path found by greedy search. & 21 & 22 \\
\bottomrule
\end{tabular}

\vspace{2pt}
{\footnotesize
All 5 certificates pass both experiments/stable\_ac/ak3\_stable\_proof/verify\_certificate.py and the independently-authored independent\_verifier.py (confirmed live: verify\_certificate.py=True, independent\_verifier.py=True).\par
}

\end{table}

\subsection{Mutation sensitivity}

Sign flips, corrupted target states, swapped intermediate states, and mis-pointed
eliminations are all rejected by both verifiers. One candidate mutation was initially
flagged as accepted; on inspection it was not a soundness hole but a genuinely valid
alternative move --- inverting a relator of length one is itself a legal AC2 move, correctly
treated as a no-op by the canonical-equality quotient rather than as a tamper. A
semantically meaningful tamper test therefore has to flip an \emph{interior} letter of a
relator of length at least three in a middle state, not merely a sign.

\subsection{Coverage and chain composition}

Of the 151 catalog leaves classified during the campaign, 14 have exported,
machine-verified leaf certificates; the remaining 137 classifications are recorded as data
(solved/unsolved, node count, floor reached) but were not individually exported as
certificates. Certificates concatenate whenever one certificate's end state matches the
next certificate's start state exactly, so a full stable-AC chain --- stabilize $\to$
substitution path $\to$ eliminate $\to$ 2-generator trivialization path, further composed
with the P25 bridge of \ref{app:misprint} --- assembles automatically into a single
machine-checkable certificate. This composition machinery already produces every
certificate reported above and would apply unchanged to any future solve.
